\documentclass[10pt,a4paper]{amsart}
\usepackage[margin=19mm]{geometry}
\usepackage{amsmath,amssymb,mathtools}
\usepackage[scr=rsfs,cal=euler]{mathalfa}
\usepackage{xfrac}
\usepackage[unicode,hidelinks,bookmarks=false]{hyperref}
\newcommand{\ie}{i.e.\ }
\newcommand{\cf}{cf.\ }
\newcommand{\resp}{respectively\ }
\newcommand{\Subsection}{\subsection}
\providecommand{\nobreakdash}{\nobreak\hbox{-}}
\setlength{\emergencystretch}{2em}
\multlinegap0pt
\usepackage{amssymb}
\usepackage{mathtools}
\usepackage[shortlabels]{enumitem}
\usepackage{xcolor}
\newtheorem{theorem}{Theorem}
\title{Fractional very fast diffusion equations in Lebesgue spaces: uniqueness and smoothing effects}
\author{Mohammed-El-Mahdi Boudaoud, Arturo de Pablo, Fernando Quir\'os}
\date{Source: arXiv:2609.12703v1 (CC BY 4.0)}
\begin{document}
\maketitle

\noindent\emph{Source and licence.} Fractional very fast diffusion equations in Lebesgue spaces: uniqueness and smoothing effects. Version arXiv:2609.12703v1 (CC BY 4.0), distributed by arXiv under the Creative Commons Attribution 4.0 International licence, http://creativecommons.org/licenses/by/4.0/, as recorded in the arXiv licence metadata for this exact version and shown on the abstract page https://arxiv.org/abs/2609.12703v1. Full licence notice: \texttt{licenses/arxiv-2609.12703-CC-BY-4.0.txt}.\par\smallskip

\noindent\textit{Editorial scope.} Counted unit: the article's theorem continuity (source0.tex lines 876--878). The article's complete original proof of it is delivered with it (source0.tex lines 880--960, one \texttt{proof} environment of 81 lines). No mathematics is written, repaired or re-derived: the statement and the proof are the article's own lines, in the article's own notation, and no retained line is substituted or re-flowed.  Retained uncounted as the source-local context the proof needs: lines 87--92.  Nothing was omitted from any retained span, so the package carries no omission record.  No retained line contains a citation, so the wrapper carries no bibliography and no bibliography entry.  No retained line contains a figure environment, an image-inclusion command or drawing code, so no image is delivered and the package declares no asset dependency.  Editorial formatting only, in addition: an AMS \texttt{amsart} preamble that restates the article's own declarations and macros used by the retained lines, namely the environments \texttt{theorem} on separate counters, together with the standard packages those lines need.  The retained environments print in this document as Theorem 1 in order of appearance, whereas the article prints the same items with the numbering of its own full document.  The complete native source of the article as distributed in the arXiv e-print for this exact version is bundled unchanged as \texttt{sources/210141/source0.tex}.
\begin{equation}\label{eq:main}
    \begin{cases}
        \partial_t u+(-\Delta)^{\frac\sigma2}|u|^{m-1}u=0,&x\in \mathbb{R}^N,\ t>0,\\
        u(x,0)=u_0(x),&x\in\mathbb{R}^N,
    \end{cases}
\end{equation}

\begin{theorem}[Continuity in $L^p$]\label{Lp continuity}
    Let $u_0 \in L^p(\mathbb{R}^N) $ with $p > p^*$. Then the corresponding $L^p$ solution \( u \) to~\eqref{eq:main} satisfies $u \in C((0,\infty); L^p(\mathbb{R}^N))$.
\end{theorem}

\begin{proof}
We will prove that for every $\varepsilon>0$ there exists $h>0$ small such that
$$
    \int_{\mathbb{R}^N}|u(x,t_1)-u(x,t_2)|^p\,dx< c(p)\varepsilon\qquad \textup{if }|t_1-t_2|<h,\ t_1,t_2>0,
$$
so in fact we have uniform continuity. We split the integral to be estimated as
$$
    \int_{\mathbb{R}^N}|u(x,t_1)-u(x,t_2)|^p\,dx=\underbrace{\int_{B_R}|u(x,t_1)-u(x,t_2)|^p\,dx}_{I_1}+\underbrace{\int_{B^c_R}|u(x,t_1)-u(x,t_2)|^p\,dx}_{I_2}
$$
for some $R>0$ to be chosen later. Take $0<\tau<\min\{t_1,t_2\}$ a time such that $u(\tau)\in L^p(\mathbb{R}^N)$. Thus there exists some large value $R_1=R_1(\varepsilon,\tau)$ such that 
$$
    \int_{B_R^c}u^p(x,\tau)\,dx<\varepsilon\quad\textup{for every }R>R_1.
$$
We claim that  there also exists $R_2=R_2(\varepsilon,\tau)>R_1$ such that
$$
    \int_{B^c_R}u^p(x,s)\,dx<2\varepsilon\quad\textup{for every }R>R_2\textup{ and almost every }s>\tau. 
$$
This would show that $I_2\le C_p \varepsilon$ for almost every $s>\tau$. 

To prove the claim, we estimate $\displaystyle\int_{\mathbb{R}^N}\phi_Ru^p(s)$, where $\phi_R(x)=\phi(x/R)$ and $\phi$ is a smooth, nonnegative, radially nonincreasing function such that
\[
    \phi(x)=
    \begin{cases}
        1, & |x|>1,\\
        0, & |x|<1/2.
    \end{cases}
\]
Multiplying the equation by $p\phi_R u^{p-1}(s)$, and using the formula for the fractional Laplacian of a product, we get 
$$    \partial_t\int_{\mathbb{R}^N}\phi_R(x)u^p(x,s)\,dx
    =-p\int_{\mathbb{R}^N} u^m(x,s)(-\Delta)^{\frac\sigma2}(u^{p-1}(x,s)\phi_R(x))\,dx=J^1+J^2,
$$
where
\begin{align*}
    J^1&
    =-p\left(\int_{\mathbb{R}^N} u^{m+p-1}(x,s)(-\Delta)^{\frac\sigma2}\phi_R(x)\,dx+\int_{\mathbb{R}^N}\phi_R(x)u^m(x,s)(-\Delta)^{\frac\sigma2}u^{p-1}(x,s)\,dx\right),
    \\
    J^2&=c\int_{\mathbb{R}^N}u^m(x,s)\mathcal{B}(u^{p-1},\phi_R)(x,s)\,dx,\quad
    \mathcal{B}(f,g)(x)=\int_{\mathbb{R}^N}\frac{(f(x)-f(y))(g(x)-g(y))}{|x-y|^{N+\sigma}}\,dy.
\end{align*}
Using Kato’s inequality,
$$
    (-\Delta)^{\frac\sigma2}u^{m+p-1}\le \frac{p+m-1}{p-1} u^m(-\Delta)^{\frac\sigma2}u^{p-1},
$$
we can combine the two integrals defining $J^1$ into just one, obtaining
\begin{align*}
    J^1&\le\int_{\mathbb{R}^N} u^{m+p-1}(x,s)|(-\Delta)^{\frac\sigma2}\phi_R(x)|\,dx 
    \le c\|u(s)\|_p^{m+p-1}\|(-\Delta)^{\frac\sigma2}\phi_R\|_{\frac p{1-m}}
    \\&
    \le c\|u_0\|_p^{m+p-1}R^{-\sigma+\frac{N(1-m)}{p}}\|(-\Delta)^{\frac\sigma2}\phi\|_{\frac p{1-m}}\le cR^{-\sigma+\frac{N(1-m)}{p}}.
\end{align*}
On the other hand,  we simplify the integral involving $\mathcal{B}$ using the estimate
$$
    u^m(x,s)u^{p-1}(y,s)\le c \left(u^{m+p-1}(x,s)+u^{m+p-1}(y,s)\right)
$$
which follows from Young’s inequality $a^p b^q \le \frac{p}{p+q}\, a^{p+q} + \frac{q}{p+q}\, b^{p+q}$, and symmetry. Thus,
\begin{align*}
J^2&\le c\int_{\mathbb{R}^N}u^m(x,s)\mathcal{B}(u^{p-1},\phi_R)(x,s)\,dx \\ &\le c
\int_{\mathbb{R}^N}u^{m+p-1}(x,s)\left(\int_{\mathbb{R}^N}\frac{|\phi_R(x)-\phi_R(y)|}{|x-y|^{N+\sigma}}\,dy\right)\,dx
\\&
 \leq cR^{-\sigma}\|u(t)\|_p^{m+p-1}\left(\int_{\mathbb{R}^N}\left( \int_{\mathbb{R}^N}\frac{|\phi_R(x)-\phi_R(y)|}{|x-y|^{N+\sigma}}\,dy\right)^{\frac{1-m}{p}}\,dx\right)^{\frac p{1-m}}
 \\&
\le  cR^{-\sigma+\frac{N(1-m)}{p}}\left(\int_{\mathbb{R}^N}\left( \int_{\mathbb{R}^N}\frac{|\phi(z)-\phi(w)|}{|z-w|^{N+\sigma}}\,dw\right)^{\frac{1-m}{p}}\,dz\right)^{\frac p{1-m}}
 \\&
\le  cR^{-\sigma+\frac{N(1-m)}{p}}.
\end{align*}
We therefore get,
$$
\partial_t\int_{\mathbb{R}^N}\phi_R(x)u^p(x,s)\,dx\le cR^{-\sigma+\frac{N(1-m)}{p}},
$$
which integrated over the time interval $(\tau,s)$  gives
$$
\int_{\mathbb{R}^N}\phi_R(x)u^p(x,s)\,dx\le \int_{\mathbb{R}^N}\phi_R(x)u^p(x,\tau)\,dx+cR^{-\sigma+\frac{N(1-m)}{p}}. 
$$
Since $p>p^*$, the exponent $-\sigma+\frac{N(1-m)}{p}$ is negative, and the claim follows if $R$ is big enough. 

Now we turn to $I_1$, which is easier to estimate. Since $p>p^*$, $u$ is bounded, and we get
$$
    I_1=\int_{B_R}|u(x,t_1)-u(x,t_2)|^p\,dx \le c(\tau,R)\int_{B_R}|u(x,t_1)-u(x,t_2)|\,dx<\varepsilon
$$
for any fixed $R$, if $|t_1-t_2|<h$ with $h$ small,  by the continuity in $L^1_{\text{loc}}$. 
\end{proof}

\end{document}
